% ============================================================
%  part1/ch04-memory.tex
%  Chapter 4: Memory
% ============================================================

\chapter{Memory}
\label{ch:memory}

\epigraph{%
  Memory is not an instrument for surveying the past\\
  but its theatre.%
}{Walter Benjamin}

\section{The reconstruction that feels like recall}

Human memory is the most familiar example of reconstruction.

Nobody stores every sensory detail of an experience.
The sights, sounds, smells, and physical sensations of an event
are not recorded like a video file.
Instead, traces are laid down --- partial, compressed,
context-dependent records --- and later experience involves
reconstruction from those traces, not playback of a recording.

Psychologists have known this since the work of Frederic Bartlett
in the 1930s.
Memory is reconstructive.
It is influenced by prior knowledge, subsequent experience,
and the context in which recall takes place.

The mathematical structure of this process is surprisingly clean.

\section{Encoding, trace, and reconstruction}

Let $x \in X$ be an experienced event, where $X$ is the space of
possible experiences.
Encoding is a map
\[
  F : X \to T,
\]
where $T$ is the space of memory traces.

Recall is a reconstruction map
\[
  G : T \to \hat{X},
\]
where $\hat{X}$ is the space of recalled experiences.

The remembered event is
\[
  \hat{x} = G(F(x)).
\]

This is the first explicit appearance of the composition
$G \circ F$ that will become the central operator of the book.

\begin{definition}[Reconstruction operator]
  The \emph{reconstruction operator} is the composition
  \[
    R = G \circ F : X \to \hat{X}.
  \]
  Memory is exact if $R(x) = x$ for all $x \in X$.
  Memory is admissible if $R(x) \approx x$ in a sense made
  precise by the reconstruction error.
\end{definition}

\section{Reconstruction error}

Memory is exact only if
\[
  G(F(x)) = x.
\]
Usually,
\[
  G(F(x)) \neq x,
\]
but recall may still be admissible if the error is small:
\[
  d(x, G(F(x))) < \varepsilon.
\]

\begin{definition}[Reconstruction error]
  The \emph{reconstruction error} of the memory system at
  experience $x$ is
  \[
    E(x) = d(x, G(F(x))),
  \]
  where $d$ is a metric on the space of experiences.
\end{definition}

This is the central quantity of CLIO (Part~V) and the
admissibility framework throughout the book.
It first appears here as a measure of how faithfully a memory
system reconstructs experience.

A memory system is:
\begin{itemize}[leftmargin=2em]
  \item \emph{exact} if $E(x) = 0$ for all $x$;
  \item \emph{admissible} if $E(x) < \varepsilon$ for some
    tolerance $\varepsilon > 0$ appropriate to the task;
  \item \emph{inadmissible} if $E(x) \geq \varepsilon$.
\end{itemize}

\section{What memory preserves}

If memory is not exact, what does it preserve?

Bartlett's original experiments showed that memory tends to
preserve:
\begin{itemize}[leftmargin=2em]
  \item overall structure and narrative;
  \item emotionally significant details;
  \item elements consistent with prior knowledge and expectations.
\end{itemize}

Memory tends to lose:
\begin{itemize}[leftmargin=2em]
  \item precise numerical detail;
  \item peripheral sensory information;
  \item elements inconsistent with prior knowledge (which are
    often systematically distorted toward the expected).
\end{itemize}

In the language developed so far, memory is a projection whose
kernel contains sensory precision and peripheral detail, while
preserving semantic structure and emotional salience.

\begin{definition}[Semantic admissibility]
  A memory reconstruction $\hat{x} = G(F(x))$ is
  \emph{semantically admissible} if it preserves the
  structural and relational features of $x$ that are relevant
  to subsequent reasoning and action, even if precise sensory
  detail is lost.
\end{definition}

This anticipates the formal admissibility conditions of Part~V,
where admissibility is defined not as exact reconstruction but
as preservation of the structure relevant to a given task or
reconstruction goal.

\section{The encoding map and its kernel}

The encoding map $F : X \to T$ necessarily has a nontrivial kernel.

\begin{proposition}[Memory compression]
  If the space of memory traces $T$ has lower dimension or
  capacity than the space of experiences $X$, then $F$ is
  non-injective: there exist $x_1 \neq x_2$ with
  $F(x_1) = F(x_2)$.
\end{proposition}

Two experiences that produce the same memory trace cannot be
distinguished by subsequent recall.

\begin{example}[Source monitoring failure]
  A well-documented phenomenon in cognitive psychology is
  \emph{source monitoring failure}: remembering a fact without
  remembering where it was learned.
  The content is encoded but the source tag is lost.
  Two routes to the same fact --- reading it versus being told it
  --- can produce the same memory trace.
  This is a concrete instance of the kernel of $F$ containing
  source information.
\end{example}

\section{Memory as the first reconstruction system}

The memory chapter introduces three ideas that will recur
throughout the book.

\begin{enumerate}[leftmargin=2em]
  \item \textbf{The reconstruction operator} $G \circ F$ is the
    natural mathematical object for studying observation-dependent
    knowledge.

  \item \textbf{The reconstruction error} $E(x) = d(x, G(F(x)))$
    measures the fidelity of reconstruction and will become the
    central quantity in CLIO and the admissibility framework.

  \item \textbf{Admissibility without exactness.}
    A reconstruction need not be exact to be useful.
    It must preserve the structure relevant to subsequent
    inference and action.
\end{enumerate}

\principle{
  Reconstruction need not be exact.\\[4pt]
  It must be admissible: sufficiently faithful\\[4pt]
  for the purposes that reconstruction serves.
}
